\documentclass[11pt]{article}

\usepackage{amsmath,amssymb}

\usepackage[margin=24mm]{geometry}

\usepackage{longtable}

\allowdisplaybreaks

\newcommand{\Sfun}{\mathcal{S}}

\newcommand{\Cfun}{\mathcal{C}}

\newcommand{\Tfun}{\mathcal{T}}

\title{SoftArm EXTENSIBLE\_EULER\_BERNOULLI + PCS Model (2 Sections)}

\author{Generated by SoftArm Toolbox}

\date{}

\begin{document}

\maketitle

\section{Model Summary}
This document is generated from the same SymPy model used for numerical code generation.
\begin{description}
\item[Source configuration] \texttt{extensible\_euler\_bernoulli\_pcs\_three\_tendon\_n2.toml}
\item[Rod theory] \texttt{extensible\_euler\_bernoulli}
\item[Spatial parameterization] \texttt{pcs}
\item[Number of sections] 2
\item[Base mode] \texttt{fixed}
\item[Arm inertia model] \texttt{lumped}
\item[Material integration] \texttt{analytic}
\end{description}

\section{Notation and Parameters}
The arm base is fixed, hence $\boldsymbol q_B$ is empty.
\begin{equation}
\boldsymbol q_a=\begin{bmatrix}b_{x,1}\\b_{y,1}\\L_{1}\\b_{x,2}\\b_{y,2}\\L_{2}\end{bmatrix}
\end{equation}

\begin{equation}
\dot{\boldsymbol q}_a=\begin{bmatrix}\dot{b_{x,1}}\\\dot{b_{y,1}}\\\dot{L_{1}}\\\dot{b_{x,2}}\\\dot{b_{y,2}}\\\dot{L_{2}}\end{bmatrix}
\end{equation}

\begin{equation}
\boldsymbol q=\begin{bmatrix}\boldsymbol q_B\\\boldsymbol q_a\end{bmatrix},\qquad \dot{\boldsymbol q}=\begin{bmatrix}\dot{\boldsymbol q}_B\\\dot{\boldsymbol q}_a\end{bmatrix}
\end{equation}

The NED world frame is used; gravity acts in the positive world $z$ direction.
\subsection{Runtime Parameters}
\begin{longtable}{lll}
\hline
Symbol & Code name & Default \\
\hline
\endfirsthead
\hline Symbol & Code name & Default \\
\hline
\endhead
$L_{0,1}$ & \texttt{s1\_rest\_length} & $0.45$ \\
$L_{0,2}$ & \texttt{s2\_rest\_length} & $0.5$ \\
$m_{1}$ & \texttt{s1\_mass} & $0.18$ \\
$m_{2}$ & \texttt{s2\_mass} & $0.2$ \\
$I_{xx,1}$ & \texttt{s1\_Ixx} & $0.0018$ \\
$I_{xx,2}$ & \texttt{s2\_Ixx} & $0.002$ \\
$I_{yy,1}$ & \texttt{s1\_Iyy} & $0.0018$ \\
$I_{yy,2}$ & \texttt{s2\_Iyy} & $0.002$ \\
$I_{zz,1}$ & \texttt{s1\_Izz} & $0.0001$ \\
$I_{zz,2}$ & \texttt{s2\_Izz} & $0.00011$ \\
$k_{b_x,1}$ & \texttt{s1\_k\_bx} & $1.1$ \\
$k_{b_x,2}$ & \texttt{s2\_k\_bx} & $1.2$ \\
$k_{b_y,1}$ & \texttt{s1\_k\_by} & $1.1$ \\
$k_{b_y,2}$ & \texttt{s2\_k\_by} & $1.2$ \\
$k_{L,1}$ & \texttt{s1\_k\_l} & $45$ \\
$k_{L,2}$ & \texttt{s2\_k\_l} & $50$ \\
$d_{b_x,1}$ & \texttt{s1\_d\_bx} & $0.05$ \\
$d_{b_x,2}$ & \texttt{s2\_d\_bx} & $0.05$ \\
$d_{b_y,1}$ & \texttt{s1\_d\_by} & $0.05$ \\
$d_{b_y,2}$ & \texttt{s2\_d\_by} & $0.05$ \\
$d_{L,1}$ & \texttt{s1\_d\_l} & $0.1$ \\
$d_{L,2}$ & \texttt{s2\_d\_l} & $0.1$ \\
$g$ & \texttt{gravity} & $9.81$ \\
$m_e$ & \texttt{tip\_mass} & $0.1$ \\
$I_{e,xx}$ & \texttt{tip\_Ixx} & $0.01$ \\
$I_{e,yy}$ & \texttt{tip\_Iyy} & $0.01$ \\
$I_{e,zz}$ & \texttt{tip\_Izz} & $0.01$ \\
$r_{\mathrm{t1},1}$ & \texttt{act\_t1\_s1\_radius} & $0.0275$ \\
$r_{\mathrm{t1},2}$ & \texttt{act\_t1\_s2\_radius} & $0.03$ \\
$r_{\mathrm{t2},1}$ & \texttt{act\_t2\_s1\_radius} & $0.0275$ \\
$r_{\mathrm{t2},2}$ & \texttt{act\_t2\_s2\_radius} & $0.03$ \\
$r_{\mathrm{t3},1}$ & \texttt{act\_t3\_s1\_radius} & $0.0275$ \\
$r_{\mathrm{t3},2}$ & \texttt{act\_t3\_s2\_radius} & $0.03$ \\
\hline
\end{longtable}

\section{Kinematics}
\begin{equation}
H_{WB}=I_4
\end{equation}

\begin{equation}
H_{BA}=\begin{bmatrix}1 & 0 & 0 & 0.0\\0 & 1 & 0 & 0.0\\0 & 0 & 1 & 0.0\\0 & 0 & 0 & 1\end{bmatrix}\label{eq:mount-transform}
\end{equation}

\subsection{Piecewise-Constant-Strain Parameterization}
\begin{equation}
\Sfun(z)=\begin{cases}\dfrac{\sin\sqrt z}{\sqrt z},&z\ne0\\1,&z=0\end{cases},\qquad \Cfun(z)=\begin{cases}\dfrac{1-\cos\sqrt z}{z},&z\ne0\\\dfrac12,&z=0\end{cases}
\end{equation}

\begin{equation}
\Tfun(z)=\begin{cases}\dfrac{1-\Sfun(z)}{z},&z\ne0\\\dfrac16,&z=0\end{cases}
\end{equation}

\begin{equation}
\kappa_i=[-b_{y,i}/L_{0,i},\ b_{x,i}/L_{0,i},\ 0]^T,\qquad\nu_i=[0,0,l_i/L_{0,i}]^T
\end{equation}

\begin{equation}
\Omega_i(\xi)=L_i\xi\widehat{\kappa_i},\qquad z_i(\xi)=(L_i\xi)^2\kappa_i^T\kappa_i
\end{equation}

\begin{equation}
R_i=I_3+\Sfun(z_i)\Omega_i+\Cfun(z_i)\Omega_i^2,\qquad r_i=\left[I_3+\Cfun(z_i)\Omega_i+\Tfun(z_i)\Omega_i^2\right]L_i\xi\nu_i
\end{equation}

\begin{equation}
H_i(\xi)=\begin{bmatrix}R_i(\xi)&r_i(\xi)\\0&1\end{bmatrix}\label{eq:local-transform}
\end{equation}

\begin{equation}
H_{Wi}(\xi)=H_{WB}H_{BA}\left(\prod_{k=1}^{i-1}H_k(1)\right)H_i(\xi)
\end{equation}

\begin{equation}
J_{v,i}=\frac{\partial r_{Wi}}{\partial\boldsymbol q},\qquad J_{\omega,i}^{(:,j)}=\operatorname{vex}\!\left(\operatorname{skew}\!\left(\frac{\partial R_{Wi}}{\partial q_j}R_{Wi}^{T}\right)\right)
\end{equation}

\begin{equation}
J_e=\begin{bmatrix}J_{v,e}\\J_{\omega,e}\end{bmatrix}
\end{equation}


\section{Energy and Dynamics}
For the lumped inertia option, each section contribution is evaluated at $\xi=1/2$.
\begin{equation}
M_i=\int_0^1\!\left(m_iJ_{v,i}^{T}J_{v,i}+J_{\omega,i}^{T}R_{Wi}I_iR_{Wi}^{T}J_{\omega,i}\right)d\xi
\end{equation}

\begin{equation}
M_e=m_eJ_{v,e}^{T}J_{v,e}+J_{\omega,e}^{T}R_eI_eR_e^TJ_{\omega,e}
\end{equation}

\begin{equation}
M=\sum_{i=1}^{N}M_i+M_e\label{eq:mass-assembly}
\end{equation}

\begin{equation}
T=\frac12\dot{\boldsymbol q}^{T}M(\boldsymbol q)\dot{\boldsymbol q}
\end{equation}

\begin{equation}
V=V_g+V_{\mathrm{elastic}},\qquad V_g=-g\!\left(\sum_{i=1}^{N}m_i z_i(1/2)+m_ez_e\right)
\end{equation}

\begin{equation}
V_{\mathrm{elastic}}=\frac12\sum_{i=1}^{N}\left(k_{b_x,i}b_{x,i}^{2}+k_{b_y,i}b_{y,i}^{2}+k_{L,i}(L_i-L_{0,i})^2\right)
\end{equation}

\begin{equation}
D=\operatorname{diag}\!\left(\begin{bmatrix}d_{b_x,1}\\d_{b_y,1}\\d_{L,1}\\d_{b_x,2}\\d_{b_y,2}\\d_{L,2}\end{bmatrix}\right)
\end{equation}

\begin{equation}
c(\boldsymbol q,\dot{\boldsymbol q})=\frac{\partial(M\dot{\boldsymbol q})}{\partial\boldsymbol q}\dot{\boldsymbol q}-\frac12\left(\frac{\partial(\dot{\boldsymbol q}^{T}M\dot{\boldsymbol q})}{\partial\boldsymbol q}\right)^T
\end{equation}

\begin{equation}
h=c+\nabla_{\boldsymbol q}V+D\dot{\boldsymbol q}
\end{equation}

The vehicle wrench $w_B$ is expressed in the vehicle body frame, whereas the end wrench $w_e$ is expressed in the NED world frame.
\begin{equation}
Q=S_a\tau_a+B_v(\boldsymbol q)w_B+J_e^Tw_e,\qquad B_v=J_B^T\operatorname{diag}(R_{WB},R_{WB})
\end{equation}

\begin{equation}
M\ddot{\boldsymbol q}+h=Q\label{eq:dynamics}
\end{equation}

\begin{equation}
\dot x=\begin{bmatrix}\dot{\boldsymbol q}\\M^{-1}(Q-h)\end{bmatrix},\qquad x=\begin{bmatrix}\boldsymbol q\\\dot{\boldsymbol q}\end{bmatrix}
\end{equation}


\section{Actuation}
The configured actuator family is \texttt{tendon} with 3 ordered channels.
\begin{center}
\begin{tabular}{ll}
\hline
Channel & Type \\
\hline
\texttt{t1} & \texttt{unilateral} \\
\texttt{t2} & \texttt{unilateral} \\
\texttt{t3} & \texttt{unilateral} \\
\hline
\end{tabular}
\end{center}
\begin{equation}
y(\boldsymbol q_a)=\begin{bmatrix}- r_{\mathrm{t1},1} b_{x,1} - r_{\mathrm{t1},2} b_{x,2} + L_{1} + L_{2}\\- r_{\mathrm{t2},1} \left(- 0.49999999999999983 b_{x,1} + 0.86602540378443874 b_{y,1}\right) - r_{\mathrm{t2},2} \left(- 0.49999999999999983 b_{x,2} + 0.86602540378443874 b_{y,2}\right) + L_{1} + L_{2}\\- r_{\mathrm{t3},1} \left(- 0.50000000000000042 b_{x,1} - 0.8660254037844384 b_{y,1}\right) - r_{\mathrm{t3},2} \left(- 0.50000000000000042 b_{x,2} - 0.8660254037844384 b_{y,2}\right) + L_{1} + L_{2}\end{bmatrix}
\end{equation}

\begin{equation}
J_a=\frac{\partial y}{\partial\boldsymbol q_a},\qquad \gamma_a=\dot J_a\dot{\boldsymbol q}_a
\end{equation}

The actuator coordinates are tendon lengths, and positive $T$ denotes tendon tension.
\begin{equation}
\tau_a=-J_a^TT
\end{equation}

Strict actuator-coordinate acceleration is enforced with the full-system Jacobian $J_{a,f}=[0\;J_a]$.
\begin{equation}
\begin{bmatrix}M&J_{a,f}^{T}\\J_{a,f}&0\end{bmatrix}\begin{bmatrix}\ddot{\boldsymbol q}\\T\end{bmatrix}=\begin{bmatrix}Q-h\\\ddot y_{\mathrm{cmd}}-\gamma_a\end{bmatrix}
\end{equation}


\end{document}
